\documentclass[11pt]{article}
\usepackage[a4paper,margin=25mm]{geometry}
\usepackage{amsmath,amssymb}
\usepackage[hidelinks]{hyperref}
\setlength{\emergencystretch}{3em}
\title{Conjecture 00000004105:\\incompatible growth clauses at an actual star}
\author{lizaixi01, with OpenAI Codex assistance}
\date{23 October 2026}
\begin{document}
\maketitle
\section{The necessary conjunction}
The original bilingual source defines a motion group using a graph's
configuration space, then claims that its growth is polynomial if and
only if the graph is a tree. The same conjecture describes the star-graph
motion group as strictly super-polynomial (and sub-exponential).
These two growth clauses cannot hold simultaneously, regardless of
which candidate growth functions are assigned to the graphs.

We prove this universal obstruction, rather than assigning a value to
the growth of a configuration-space fundamental group. The particle
number and configuration-space conventions are unspecified in the
source; the argument applies to any consistent choice giving a growth
function for each relevant graph. No fundamental group or its actual
growth is computed here. The sub-exponential and minimality assertions
are unnecessary to the contradiction.

\section{Actual graph and quantitative growth}
Let $S$ be the actual four-vertex star, with center $0$ and vertex set
$\{0,1,2,3\}$. In Lean it is Mathlib's
\texttt{SimpleGraph.starGraph (0 : Fin 4)}. The library's genuine
\texttt{IsTree} theorem proves that $S$ is a tree; its actual edge set
has cardinality three.

For an arbitrary candidate $f:\mathbb N\to\mathbb R$, polynomial
upper growth means
\[
 \exists d\in\mathbb N,\qquad f(n)=O(n^d)\quad(n\to\infty).
\]
The code uses Mathlib's actual \texttt{IsBigO} and proves its equivalence to
\[
 \exists d,\ C>0,\ N,\quad
 \forall n\geq N:\quad |f(n)|\leq Cn^d.                     \tag{1}
\]
The quantitative super-polynomial condition is
\[
 \forall d\in\mathbb N\ \forall C>0\ \forall N\in\mathbb N,
 \quad \exists n\geq N:\quad Cn^d<f(n).                    \tag{2}
\]
It is defined through numerical values and thresholds, not as a
syntactic negation. For nonnegative count functions, the proof also
shows that (2) is equivalent to the absence of every polynomial
big-$O$ upper bound. It does not demand the stronger condition of
eventual domination at every subsequent index.

\section{The contradiction and its scope}
Suppose a candidate graph-to-growth assignment satisfies the source's
tree-to-polynomial implication and the star's asserted (2).
Since $S$ is a tree, its candidate function $f_S$ satisfies (1), with
some actual $d,C,N$. Applying (2) to those same values gives $n\geq N$
with
\[
 Cn^d<f_S(n)\leq |f_S(n)|\leq Cn^d,
\]
a contradiction. Both thresholds and constants belong to the same
function; no numerical growth value is substituted.

The proof also covers every finite vertex type and every star center,
so it does not depend on the star's unspecified order. The final
Lean theorems quantify over every assignment to actual
four-vertex graphs, then explicitly over assignments to all finite
labelled graphs, including natural-valued candidate counting functions.
Thus any actual motion-group growth assignment is covered if it is to
satisfy both source clauses. This does not prove any standalone
classification of graph configuration spaces, any triviality of tree
motion groups, or any corrected conjecture omitting either clause.

\section{Prior failed submission and verification}
The contradiction idea was already public in
\href{https://github.com/The-Last-Math-Competition/The-Last-Math-Competition/pull/293}{PR 293}.
Its \href{https://github.com/The-Last-Math-Competition/The-Last-Math-Competition/pull/293\#issuecomment-5968931621}{maintainer review}
rejected a Lean theorem consisting only of fixed natural-number
arithmetic, without actual graph or growth objects. The present proof
repairs those omissions using a real star/tree, arbitrary growth
functions, actual big-$O$, and all numerical constants and cutoffs.
We do not inherit the prior PR's unsupported claim that tree motion
groups are trivial.

All named local declarations are audited in a pinned Lean 4.33.0 and
Mathlib project. The package includes a new Main build, full official
fresh kernel replay, source-bound command and log receipts, original
bilingual source, LaTeX and this PDF. No \texttt{sorry},
\texttt{native\_decide} or additional axiom is used. OpenAI Codex
assisted with a separate supervising AI reviewing the proof, source
scope and rendered pages; no human certification or organizer
acceptance is claimed.
\end{document}
